\section{Model}\label{sec:model}

There are two jurisdictions, $A$ and $B$, two ex ante symmetric firms, and a representative consumer. Jurisdictions choose place-based climate protection, firms choose primary production locations and geographic backup readiness, localized disasters are realized, and surviving firms compete in the product market. The model is deliberately static. Its purpose is to isolate one interaction: local competition for mobile production changes public protection, while public protection changes firms' private production redundancy.

\subsection{Product market}

Consumer utility is
\begin{equation}
U(x_1,x_2)=x_1+x_2-\frac12\left(x_1^2+x_2^2+2\gamma x_1x_2\right),
\qquad 0\leq \gamma<1,
\label{eq:utility}
\end{equation}
where $x_i$ is consumption of firm $i$'s product and $\gamma$ indexes product substitutability. The corresponding inverse demands are
\begin{equation}
p_i=1-x_i-\gamma x_j.
\label{eq:inverse-demand}
\end{equation}
Marginal production cost is normalized to zero. Whenever a firm is available through either its primary plant or its backup arrangement, it can supply the unconstrained Cournot quantity at this same marginal cost; there is no separate state-dependent capacity constraint. If both firms are available after the disaster state is realized, they compete in quantities. If only one firm is available, it is a monopolist. If neither firm is available, output is zero.

\subsection{Public protection and localized risk}

Jurisdiction $j\in\{A,B\}$ chooses public climate protection
\begin{equation}
a_j\in[0,\bar a].
\end{equation}
Protection lowers the probability that primary production located in $j$ is disrupted:
\begin{equation}
q_j=q_0-a_j,
\qquad 0<\bar a<q_0<1.
\label{eq:risk}
\end{equation}
The resource cost of protection is $c a_j^2/2$, with $c>0$. The linear risk technology defines the baseline theorem environment. Section~\ref{sec:discussion} reports a diagnostic portability check under a decreasing nonlinear risk technology; that check is robustness evidence rather than an extension of the exact theorem to every decreasing risk function.

Localized shocks generate correlation in primary production failures. Let $F_i\in\{0,1\}$ indicate failure of firm $i$'s primary production. If both firms place their primary plants in the same jurisdiction $j$, their primary plants fail jointly with probability $J=q_j$. If they locate in different jurisdictions, regional shocks are independent and joint primary failure occurs with probability
\begin{equation}
J=q_Aq_B.
\label{eq:joint-risk}
\end{equation}
Thus $\Pr(F_i=1)=s_i$ and $\Pr(F_1=F_2=1)=J$, with the location configuration determining $s_i$ and $J$.

\subsection{Primary location and geographic backup readiness}

Firm $i$ chooses a primary production location $L_i\in\{A,B\}$. Let
\begin{equation}
s_i=q_{L_i}
\end{equation}
be the failure probability of its primary production. After primary locations are chosen, firm $i$ chooses backup readiness $r_i\in[0,1]$. Conditional on primary failure, backup production is available with probability $r_i$. Maintaining readiness costs
\begin{equation}
C_i^R(r_i)=\frac{k}{2}r_i^2,
\qquad k>0,
\label{eq:backup-cost}
\end{equation}
and this cost is paid ex ante and is state independent.

The joint law of backup success is a maintained primitive. Conditional on public policies, primary locations, and the realized primary-failure vector $(F_1,F_2)$, backup-success draws are independent across firms. Whenever $F_i=1$, firm $i$'s backup is available with conditional probability $r_i$; conditional backup success does not otherwise depend on the regional shock that generated the primary-failure vector. The backup-success draws are also independent of the location-fit shocks introduced below. Equivalently, firm $i$ is unavailable with probability $s_i(1-r_i)$, while both firms are unavailable with probability
\begin{equation}
J(1-r_1)(1-r_2).
\label{eq:joint-unavailability}
\end{equation}
This conditional-independence assumption is substantive: fixing only the marginal success probabilities $r_i$ would not determine the joint production-availability states. The reduced-form variable $r_i$ should therefore be read as the success probability of the firm's geographically separate contingency arrangement after its primary site fails; any hazard exposure specific to that arrangement is already incorporated into $r_i$, and there is no additional common backup shock in the baseline model.

This reduced-form readiness choice captures costly alternative facilities, pre-arranged production capacity, or other geographically separate production options. It is smooth by construction; none of the main results depends on crossing a fixed-cost threshold.

To avoid a knife-edge winner-take-all location response, each firm receives a privately observed idiosyncratic location-fit shifter favoring $A$ relative to $B$,
\begin{equation}
\varepsilon_i\sim U[-H,H],\qquad H>0,
\label{eq:location-shock}
\end{equation}
independently across firms and independently of all regional-disaster and backup-success draws. After public policies are chosen and before the simultaneous location decisions, firm $i$ privately observes $\varepsilon_i$. Each firm knows the distribution of the rival's shifter and forms beliefs over the rival's location choice in the usual Bayesian way. We use $\varepsilon_i$ as a smoothing device for the discrete location decision, not as consumption, accounting profit, or a real resource. It therefore affects the firm's cutoff rule but is not part of baseline material social surplus. Given continuation profits, firm $i$ chooses $A$ when its deterministic payoff advantage of $A$ plus $\varepsilon_i$ is nonnegative. Appendix~\ref{app:verification} also shows that counting the selected fit payoff in welfare leaves the canonical policy ranking unchanged.

\subsection{Governments and welfare}

Let $\mathcal S(a_A,a_B)$ denote expected national product-market surplus net of private backup costs, after location and backup choices are resolved. A primary plant located in a jurisdiction generates local employment, tax-base, or immobile-factor benefits summarized by $b>0$. The baseline incidence rule assigns one half of national market surplus $\mathcal S$ to each jurisdiction. If $p$ denotes the symmetric probability that a firm chooses $A$, jurisdiction $A$'s payoff is
\begin{equation}
G_A=\frac12\mathcal S+2bp-\frac{c}{2}a_A^2,
\label{eq:government-payoff}
\end{equation}
and $B$ is symmetric. This equal-share rule is a maintained reduced-form incidence assumption, not an accounting identity derived from the product market. A coordinated planner values the same real product-market surplus and both local plant benefits:
\begin{equation}
W=\mathcal S+2b-\frac{c}{2}(a_A^2+a_B^2).
\label{eq:national-welfare}
\end{equation}
Because there are always two primary plants, $2b$ is constant at the national level. With the equal-share incidence rule, $G_A+G_B=W$. The decentralized governments therefore have a plant-attraction motive that is absent from the planner's location accounting.

\subsection{Timing and equilibrium concept}

The game proceeds as follows.
\begin{enumerate}[label=\arabic*.]
\item Jurisdictions simultaneously choose $(a_A,a_B)$.
\item Each firm privately observes its own $\varepsilon_i$; firms then simultaneously choose primary locations, forming beliefs over the rival's private draw.
\item Firms simultaneously choose backup readiness $(r_1,r_2)$ after primary locations are observed.
\item Localized primary-disaster states are realized; conditional on primary failures, backup success is realized according to the conditional-independence law above.
\item Available firms compete in quantities.
\end{enumerate}
We solve by backward induction. The location stage is a symmetric Bayesian Nash equilibrium in cutoff strategies. The analysis uses the interior continuation region certified for the stated policy domain. Upstream deviations are evaluated by re-solving the downstream location and backup subgames rather than holding on-path continuation formulas fixed.
